% Source: https://artofproblemsolving.com/community/c3h1593664_algebra_qn
% Stable identity: aops-topic-1593664-q01
Let the integers be $a<b<c<d$, so $\dfrac{1}{a}+\dfrac{1}{b}+\dfrac{1}{c}+\dfrac{1}{d}=\dfrac{19}{10}$.

\textbf{$a=1$.} Otherwise $a\ge2$ and the sum is at most $\dfrac{1}{2}+\dfrac{1}{3}+\dfrac{1}{4}+\dfrac{1}{5}=\dfrac{77}{60}<\dfrac{19}{10}$.

\textbf{$b=2$.} Otherwise $b\ge3$ and the sum is at most $1+\dfrac{1}{3}+\dfrac{1}{4}+\dfrac{1}{5}=\dfrac{107}{60}<\dfrac{19}{10}$.

\textbf{$c$ and $d$.} Now $\dfrac{1}{c}+\dfrac{1}{d}=\dfrac{19}{10}-\dfrac{3}{2}=\dfrac{2}{5}$ with $2<c<d$. Since $\dfrac{1}{c}>\dfrac{1}{5}$, we have $c<5$, so $c=3$ or $c=4$.
\begin{itemize}
\item $c=3$: $\dfrac{1}{d}=\dfrac{2}{5}-\dfrac{1}{3}=\dfrac{1}{15}$, so $d=15$.
\item $c=4$: $\dfrac{1}{d}=\dfrac{2}{5}-\dfrac{1}{4}=\dfrac{3}{20}$, not the reciprocal of an integer.
\end{itemize}
So the integers are $1,2,3,15$ (indeed $1+\dfrac{1}{2}+\dfrac{1}{3}+\dfrac{1}{15}=\dfrac{30+15+10+2}{30}=\dfrac{57}{30}=1.9$), and their sum is
\[
1+2+3+15=\boxed{21}.
\]
